A wearable chest-mounted prototype, the \emph{SomaTune Strap}, is conceived as a 40\,Hz vibrotactile device with a slow 0.1\,Hz amplitude pulse. In practical terms, this means a carrier vibration in the gamma-range is delivered locally to the torso while its intensity rises and falls in a much slower envelope. The design is intended to combine a steady entrainment-like stimulus with a long-period modulation that may be more tolerable and more physiologically legible than continuous stimulation alone.

Placed over the chest, the strap provides a direct interface to the body’s central mass, where mechanical vibration can couple into soft tissue, respiration-linked motion, and broader autonomic state. The 40\,Hz component aligns with the frequency band often discussed in relation to neural gamma entrainment, while the 0.1\,Hz pulse introduces a slow rhythmic contour that may support comfort, pacing, and sustained wear. As a prototype concept, the SomaTune Strap emphasizes localized delivery, simple hardware integration, and waveform layering rather than high-intensity stimulation.

In the RBM framework, this device serves as a compact example of how a single wearable can combine a biologically salient carrier frequency with a slower organizational rhythm. The result is not merely a vibrating accessory, but a frequency-shaped chest interface designed to explore how layered temporal patterns may modulate physiological state.
